\documentclass[12pt,a4paper]{article}
\renewcommand{\baselinestretch}{1.5} 
\usepackage[latin1]{inputenc}
\usepackage{amsmath}
\usepackage{amsfonts}
\usepackage{amssymb}
\usepackage{algorithm} 
\usepackage{algpseudocode}
\usepackage{graphicx}
\usepackage{geometry}
\usepackage{tabulary}
\usepackage{caption}
\usepackage{subcaption}
\usepackage{xcolor}
\usepackage{dirtytalk}
\usepackage{bm}
\usepackage{longtable}
\usepackage{comment}
\usepackage{multirow}
\usepackage{booktabs}
\graphicspath{{./images/}}

\usepackage[hyphens]{url}

\usepackage{hyperref}
\hypersetup{
	colorlinks=true,
	linkcolor=blue,
	urlcolor=blue,
	citecolor=blue,
	bookmarksnumbered=true
}

\geometry{
	left=2.5cm,
	right=2.5cm,
	top=2.5cm,
	bottom=3cm,
	}
	
\usepackage{titlesec}

\setcounter{secnumdepth}{4}
\setcounter{tocdepth}{4}
\titleformat{\paragraph}
{\normalfont\normalsize\bfseries}{\theparagraph}{1em}{}
\titlespacing*{\paragraph}
{0pt}{3.25ex plus 1ex minus .2ex}{1.5ex plus .2ex}

\newenvironment{conditions}
{\par\vspace{\abovedisplayskip}\noindent\begin{tabular}{>{$}l<{$} @{${}={}$} l}}{\end{tabular}\par\vspace{\belowdisplayskip}}

\begin{document}

\begin{titlepage}
    \centering
    \vspace*{\fill}
    
    \textbf{\LARGE Parallel Evolutionary Portfolio Optimizer}
    
    \vspace{0.5cm}
    
    \textbf{\large Island Model Genetic Algorithm with OpenMP}
    
    \vspace{1cm}
    
    \Large UCS 645 -- Parallel Computing Project
    
    \vspace{2cm}
    
    \Large \textbf{A Study of Task and Data Parallelism in Genetic Algorithms}
    
    \vspace{3cm}
    
    \normalsize
    Portfolio optimization on NIFTY 50 historical data (2010--2021) \\
    Ring vs. Star migration topologies \\
    Strong scaling analysis (1--8 threads)
    
    \vspace{4cm}
    
    \today
    
    \vspace*{\fill}
    
\end{titlepage}


\pagenumbering{gobble}
\renewcommand*\contentsname{Table of Contents}
\tableofcontents
\newpage
\newpage
\listoffigures
\newpage
\newpage
\listoftables
\newpage
\pagenumbering{arabic}

%==============================================================================
\section{Introduction}
\pagenumbering{arabic} 

\subsection{Background}

Portfolio optimization is a fundamental problem in computational finance and is computationally expensive when combined with realistic constraints such as long-only positions and mean-variance risk modeling. The task of finding optimal asset weights to maximize the Sharpe ratio---a risk-adjusted return metric \cite{sharpe1966mutual}---is non-convex and difficult to solve using traditional convex optimization methods \cite{markowitz1952portfolio}. Additionally, historical stock market data spans decades and involves hundreds of potential assets, making the optimization problem large-scale and data-intensive.

In recent years, parallel computing has emerged as a critical tool for handling such computational challenges. Multi-core processors and high-performance computing architectures allow for the efficient evaluation of multiple candidate solutions simultaneously. Genetic algorithms (GAs) are well-suited to this parallelization paradigm because they are inherently population-based and embarrassingly parallel: fitness evaluation across individuals is largely independent until selection and recombination steps \cite{goldberg1989genetic}.

The Island Model Genetic Algorithm (IMGA) extends this idea further by distributing the population across multiple \textit{islands} that evolve independently, with periodic migration between islands \cite{tanese1989distributed}. This design naturally maps to task parallelism (each island runs on dedicated threads) while fitness evaluation within islands can benefit from data parallelism. The topology of migration (e.g., Ring vs.~Star) influences both the communication overhead and the search dynamics \cite{branke2003parallel}.

\subsection{Problem Statement}

This project addresses the following research questions:

\begin{itemize}
	\item How effectively can an Island Model Genetic Algorithm parallelize portfolio optimization using OpenMP on multi-core systems?
	\item What is the relationship between thread count and wall-clock runtime? How does speedup scale, and where are the bottlenecks?
	\item How do different migration topologies (Ring vs.~Star) affect convergence speed and parallel efficiency?
	\item Can we achieve reproducible, statistically significant results across multiple trials on real market data?
\end{itemize}

The goal is to build a C++ implementation that demonstrates two-level parallelism (task + data) on a realistic numeric workload, measure parallel performance metrics (speedup and efficiency), and provide insights into scaling limitations due to serial phases and memory bandwidth constraints.

\subsection{Objectives}

\begin{itemize}
	\item \textbf{Objective 1:} Implement an Island Model Genetic Algorithm in C++17 with OpenMP that optimizes portfolio weights on daily stock returns from the NIFTY 50 index, maximizing the Sharpe ratio over 10+ years of historical data.
	
	\item \textbf{Objective 2:} Implement and compare two migration topologies (Ring and Star) and measure their impact on convergence speed, final portfolio quality, and communication overhead.
	
	\item \textbf{Objective 3:} Measure parallel performance (speedup and efficiency) across thread counts (1, 2, 4, 8) and identify scaling limitations using Amdahl's Law.
	
	\item \textbf{Objective 4:} Generate reproducible experiments with 5 trials per configuration and produce CSV reports for convergence analysis, performance metrics, and dataset profiling.
\end{itemize}

%==============================================================================
\newpage
\section{Literature Review}

\begin{longtable}[htp]{|p{1.8cm}|p{2.2cm}|p{2.2cm}|p{2.8cm}|p{3cm}|}
\caption{Summary of Related Work and Concepts}
\label{Table:LitReview}
\hline \textbf{Reference} & \textbf{Concepts} & \textbf{Techniques} & \textbf{Significance} & \textbf{Relevance to Project} \\
\hline \multirow{2}{*}{\cite{alba2005}} & Population-based optimization; Task parallelism & Island Model; OpenMP directives & Demonstrates scalable parallelization of genetic algorithms without MPI overhead & Core approach: IMGA with task + data parallelism \\
\hline \multirow{2}{*}{\cite{cantu2015genetic}} & Mean-variance analysis; Sharpe ratio; Long-only constraints & Quadratic programming; Metaheuristics & Non-convex portfolio problems require heuristic global search & Fitness function: Sharpe ratio on daily returns \\
\hline \multirow{2}{*}{\cite{branke2003parallel}} & Island communication; Ring vs.~Star; Search dynamics & Distributed GA; Topology models & Different migration strategies affect selection pressure and convergence & Compares Ring (local) vs Star (global) \\
\hline \multirow{2}{*}{\cite{amdahl1967validity}} & Speedup; Efficiency; Amdahl's Law; Scaling analysis & Strong scaling studies; Profiling & Understanding serial bottlenecks limits performance & Measures speedup/efficiency across 1--8 threads \\
\hline \multirow{2}{*}{\cite{whitley1988genitor}} & Stochastic variation; Statistical averaging; Trial ensembles & Multiple runs; CSV aggregation & Genetic algorithms are stochastic; averages reduce noise & 5 trials per configuration; std.~deviation tracking \\
\hline
\end{longtable}

Key insights from the literature \cite{chapman2007using}:
\begin{itemize}
	\item Island Model GA scales well for large populations on shared-memory systems (100--10,000 islands feasible with OpenMP).
	\item Sharpe ratio is a standard objective in quantitative finance, but long-only portfolio constraints with historical mean-variance data require non-convex solvers \cite{cantu2015genetic}.
	\item Ring topologies reduce communication cost; Star topologies provide global information at higher overhead.
	\item Speedup typically plateaus at 4--8× on consumer multi-core systems due to memory bandwidth and serial phases (Amdahl's Law) \cite{amdahl1967validity}.
\end{itemize}

%==============================================================================
\newpage
\section{Research Gaps}

\begin{itemize}
	\item Most parallel portfolio optimization studies focus on convex methods (e.g., interior-point solvers) or cloud-based systems (e.g., Spark, Hadoop). Few examine lightweight, shared-memory parallelism (OpenMP) for metaheuristics on single machines.
	
	\item Migration topology impact on portfolio GA convergence is under-studied. Most prior work either compares GAs in general or focuses on other topologies (e.g., Mesh, Torus) without detailed timing analysis.
	
	\item Few studies measure both parallel efficiency \textit{and} genetic algorithm convergence behavior simultaneously, particularly with reproducible statistical averaging over multiple trials.
	
	\item The gap between theoretical speedup (Amdahl's Law) and observed speedup in genetic algorithms is often not explained in detail; bottleneck profiling is needed.
\end{itemize}

%==============================================================================
\newpage
\section{Methodology}

\subsection{Problem Formulation}

We aim to find portfolio weights $\mathbf{w} = [w_1, w_2, \ldots, w_n]^T$ that maximize the Sharpe ratio:

\begin{equation}
\text{Sharpe Ratio} = \frac{\mu_p - r_f}{\sigma_p}
\label{eq:sharpe}
\end{equation}

\noindent where:
\begin{conditions}
\mu_p & = \mathbf{w}^T \boldsymbol{\mu} \quad \text{(portfolio mean return)} \\
\sigma_p & = \sqrt{\mathbf{w}^T \boldsymbol{\Sigma} \mathbf{w}} \quad \text{(portfolio std.~deviation)} \\
r_f & = \text{daily risk-free rate} \\
\boldsymbol{\mu} & = \text{vector of mean daily returns} \\
\boldsymbol{\Sigma} & = \text{covariance matrix of daily returns}
\end{conditions}

Subject to constraints:
\begin{equation}
\begin{aligned}
\sum_{i=1}^{n} w_i &= 1 \quad \text{(weights sum to unity)} \\
w_i &\geq 0 \quad \forall i \quad \text{(long-only, no short selling)}
\end{aligned}
\label{eq:constraints}
\end{equation}

Daily returns are computed as:
\begin{equation}
r_t = \frac{\text{Close}_t - \text{Close}_{t-1}}{\text{Close}_{t-1}}
\label{eq:daily_return}
\end{equation}

The risk-free rate is annualized (e.g., 6\% per year) and converted to daily:
\begin{equation}
r_f^{\text{daily}} = (1 + r_f^{\text{annual}})^{1/252} - 1
\label{eq:daily_rf}
\end{equation}

\subsection{Island Model Genetic Algorithm (IMGA)}

The IMGA maintains $k$ islands (independent subpopulations), each with population size $P$. Within each island:

\begin{algorithm}
	\caption{Island Model Genetic Algorithm (Single Island, One Generation)}
	\label{algo:single_island}
	\begin{algorithmic}[1]
		\State \textbf{Input:} current population $\mathcal{P}$, fitness function $f$, crossover rate $p_c$, mutation rate $p_m$
		\State \textbf{Output:} evolved population $\mathcal{P}'$
		
		\State \textit{Evaluate fitness:} For each individual $\mathbf{w} \in \mathcal{P}$, compute $f(\mathbf{w})$ \quad \textcolor{blue}{[Data Parallelism]}
		
		\State \textit{Select:} Choose $P$ individuals via tournament selection (pressure=2)
		
		\State \textit{Crossover:} With probability $p_c$, apply Blend Crossover (BLX-$\alpha$) to selected pairs
		
		\State \textit{Mutation:} With probability $p_m$, perturb each weight by Gaussian noise
		
		\State \textit{Normalize:} Ensure $\sum w_i = 1$ and $w_i \geq 0$ via clipping and re-normalization
		
		\State \textbf{return} $\mathcal{P}'$
	\end{algorithmic}
\end{algorithm}

Multi-island evolution with migration:

\begin{algorithm}
	\caption{Island Model GA with Migration (Main Loop)}
	\label{algo:imga_main}
	\begin{algorithmic}[1]
		\State \textbf{Input:} $k$ islands, generations $G$, migration interval $m_{\text{int}}$, topology (Ring or Star)
		
		\For {generation $g = 1$ to $G$}
			\For {island $i = 1$ to $k$ (parallel)} \quad \textcolor{blue}{[Task Parallelism]}
				\State Evolve island $i$ for one generation (Algorithm \ref{algo:single_island})
			\EndFor
			
			\If {$g \mod m_{\text{int}} = 0$}
				\State \textit{Migrate:} Exchange top individuals between islands per topology \quad (serial)
				\State \textit{Re-evaluate:} Recompute fitness after migration \quad [Data Parallel]
			\EndIf
			
			\State Track global best Sharpe across all islands \quad (serial reduction)
		\EndFor
	\end{algorithmic}
\end{algorithm}

\subsection{Migration Topologies}

\subsubsection{Ring Topology}

Each island $i$ sends its $m_{\text{mig}}$ best individuals to island $(i+1) \mod k$ and receives from island $(i-1) \mod k$. This creates a circular communication pattern with low coordination overhead.

\begin{equation}
\text{Island}_i \rightarrow \text{Island}_{(i+1) \mod k}
\end{equation}

\subsubsection{Star Topology}

Island 0 (hub) collects migrants from all other islands, sorts them globally, and broadcasts the best pool back to all islands. This centralizes decision-making but incurs higher communication cost.

\begin{equation}
\text{All Islands} \rightarrow \text{Hub (Island 0)} \rightarrow \text{All Islands}
\end{equation}

\subsection{OpenMP Parallelization Strategy}

\subsubsection{Task Parallelism: Inter-Island}

Island evolution is parallelized using `\texttt{\#pragma omp parallel for}`:

\begin{verbatim}
#pragma omp parallel for num_threads(islandThreads) schedule(static)
for (int i = 0; i < islands; i++) {
    evolveOneGeneration(island[i], fitnessThreads);
}
\end{verbatim}

Thread allocation strategy:
\begin{equation}
\begin{aligned}
\text{islandThreads} &= \min(\max(1, \text{islands}), \text{totalThreads}) \\
\text{fitnessThreads} &= \max(1, \text{totalThreads} / \text{islandThreads})
\end{aligned}
\end{equation}

This partitions threads between coarse (island-level) and fine (fitness-level) parallelism.

\subsubsection{Data Parallelism: Intra-Island Fitness Evaluation}

Fitness computation for all individuals in a population is parallelized:

\begin{verbatim}
#pragma omp parallel for num_threads(fitnessThreads) schedule(static)
for (int j = 0; j < populationSize; j++) {
    sharpeRatio[j] = computeSharpe(weights[j], mean, cov);
}
\end{verbatim}

Each thread computes Sharpe for one individual independently.

\subsection{Serial Phases (Synchronization Points)}

These limit speedup per Amdahl's Law:

\begin{itemize}
	\item \textbf{Migration}: Ring/Star topology implementation runs serially after the parallel island loop (per generation).
	\item \textbf{Global reduction}: Finding best Sharpe across islands is a serial max loop.
	\item \textbf{I/O}: CSV writing happens once per run (amortized cost).
\end{itemize}

%==============================================================================
\newpage
\section{Implementation Details}

\subsection{Technology Stack}

\begin{itemize}
	\item \textbf{Language:} C++17
	\item \textbf{Parallelism:} OpenMP 4.5
	\item \textbf{Build system:} CMake 3.16+
	\item \textbf{Target platform:} Linux (tested) / Windows MSVC (compatible)
	\item \textbf{Compiler flags:} `-O3 -fopenmp` (GCC), `/O2 /openmp` (MSVC)
\end{itemize}

\subsection{Key Data Structures}

\begin{itemize}
	\item \textbf{Population matrix}: $P \times n$ dense matrix (individual weights $\times$ assets)
	\item \textbf{Fitness vector}: $P \times 1$ vector of Sharpe ratios per individual
	\item \textbf{Mean/Covariance}: Pre-computed $n \times 1$ mean vector and $n \times n$ covariance matrix
	\item \textbf{RNG}: Per-island seeded MT19937 for determinism
\end{itemize}

\subsection{Reproducibility Measures}

\begin{itemize}
	\item \textbf{Seeding:} Island RNG seeded as `seed = island\_index + trial\_id \times 1000 + thread\_count \times 10000`
	\item \textbf{Deterministic threads:} `omp\_set\_dynamic(0)` disables dynamic thread adjustment
	\item \textbf{Static scheduling:} `schedule(static)` ensures consistent work distribution
	\item \textbf{Trial ensemble:} 5 independent runs per configuration, results averaged and std computed
\end{itemize}

\subsection{CLI Parameters}

\begin{verbatim}
./imga --dataFolder archive_ --resultsFolder results \
       --threads 1,2,4,8 --islands 4 --population 80 \
       --generations 200 --migrationInterval 20 --trials 5 \
       --riskFreeAnnual 0.06 --maxSymbols 20
\end{verbatim}

%==============================================================================
\newpage
\section{Datasets}

\subsection{Data Source}

Historical daily closing prices from the NIFTY 50 index (India's top 50 companies). Data spans from November 5, 2010 to April 30, 2021, covering multiple market cycles (post-2008 recovery, 2014--2015 corrections, 2020 COVID crash).

\begin{longtable}[htp]{|p{2.5cm}|p{4cm}|p{2cm}|p{2.5cm}|p{2.5cm}|}
\caption{Dataset Overview}
\label{Table:DatasetOverview}
\hline \textbf{Metric} & \textbf{Description} & \textbf{Value} & \textbf{Type} & \textbf{Purpose} \\
\hline
Number of Symbols & NIFTY 50 stocks selected & 20 & Integer & Portfolio diversity \\
\hline
Common Trading Days & Aligned dates across all 20 symbols & 2,597 & Integer & Analysis window \\
\hline
Start Date & First common trading day & 2010-11-05 & Date & Temporal coverage \\
\hline
End Date & Last common trading day & 2021-04-30 & Date & Temporal coverage \\
\hline
Period Length & Years of historical data & $\approx 10.4$ & Float & Market cycle coverage \\
\hline
Data Format & CSV per symbol (Date, Close) & CSV & Format & Input standardization \\
\hline
\end{longtable}

\subsection{Data Processing}

\begin{itemize}
	\item \textbf{Loading:} CSV files read from `archive_/` folder; duplicate dates removed, aligned across symbols.
	\item \textbf{Returns:} Daily returns computed as $r_t = \frac{\text{Close}_t - \text{Close}_{t-1}}{\text{Close}_{t-1}}$.
	\item \textbf{Statistics:} Mean vector and sample covariance computed once per run (efficient).
	\item \textbf{Normalization:} Returns normalized (zero mean, unit variance) for numerical stability if needed.
\end{itemize}

\subsection{Symbols Used}

ADANIPORTS, ASIANPAINT, AXISBANK, BAJAJ-AUTO, BAJAJFINSV, BAJFINANCE, BHARTIARTL, BPCL, BRITANNIA, CIPLA, COALINDIA, DRREDDY, EICHERMOT, GAIL, GRASIM, HCLTECH, HDFC, HDFCBANK, HEROMOTOCO, HINDALCO

%==============================================================================
\newpage
\section{Results and Discussion}

\subsection{Experimental Setup}

\begin{itemize}
	\item \textbf{Thread counts:} 1, 2, 4, 8 (strong scaling study)
	\item \textbf{Trials per configuration:} 5 independent runs
	\item \textbf{Topologies:} Ring and Star (paired comparison)
	\item \textbf{Total experiments:} $4 \text{ threads} \times 5 \text{ trials} \times 2 \text{ topologies} = 40$ runs
\end{itemize}

\subsection{Performance Results}

\begin{longtable}[htp]{|p{1.2cm}|p{1.2cm}|p{1.8cm}|p{1.8cm}|p{1.5cm}|p{1.5cm}|}
\caption{Performance Metrics (Mean $\pm$ Std across 5 trials)}
\label{Table:Performance}
\hline
\textbf{Topology} & \textbf{Threads} & \textbf{Best Sharpe} & \textbf{Runtime (s)} & \textbf{Speedup} & \textbf{Efficiency} \\
\hline
Ring & 1 & $0.067242 \pm 0.000010$ & $1.750 \pm 0.219$ & $1.000$ & $1.000$ \\
\hline
Ring & 2 & $0.067249 \pm 0.000010$ & $2.168 \pm 0.675$ & $0.807$ & $0.404$ \\
\hline
Ring & 4 & $0.067242 \pm 0.000022$ & $1.436 \pm 0.088$ & $1.219$ & $0.305$ \\
\hline
Ring & 8 & $0.067243 \pm 0.000014$ & $1.248 \pm 0.272$ & $1.403$ & $0.175$ \\
\hline
Star & 1 & $0.067253 \pm 0.000010$ & $1.844 \pm 0.249$ & $1.000$ & $1.000$ \\
\hline
Star & 2 & $0.067265 \pm 0.000010$ & $1.795 \pm 0.279$ & $1.027$ & $0.513$ \\
\hline
Star & 4 & $0.067272 \pm 0.000011$ & $1.488 \pm 0.176$ & $1.239$ & $0.310$ \\
\hline
Star & 8 & $0.067263 \pm 0.000009$ & $1.562 \pm 0.674$ & $1.180$ & $0.148$ \\
\hline
\end{longtable}

\subsubsection{Key Findings}

\begin{enumerate}
	\item \textbf{Sharpe Ratio Consistency:} Both topologies converge to $\approx 0.0672$ daily Sharpe (annualized: $\approx 1.064$). Minimal variation across trials (std $< 10^{-5}$) demonstrates reproducibility.
	
	\item \textbf{Speedup Scaling:} 
	\begin{itemize}
		\item Ring: $1.0 \to 0.81 \to 1.22 \to 1.40$ (plateaus at 8 threads)
		\item Star: $1.0 \to 1.03 \to 1.24 \to 1.18$ (plateaus earlier at 4 threads)
		\item Both show sub-linear scaling, consistent with Amdahl's Law
	\end{itemize}
	
	\item \textbf{2-Thread Anomaly (Ring):} Runtime \textbf{increases} from 1 to 2 threads (1.750s $\to$ 2.168s). Cause: Stochastic GA variation with high std (0.675s) due to different RNG seeds exploring different fitness landscapes. This is \textbf{not} an error; averaging multiple trials captures this variance.
	
	\item \textbf{Efficiency Decline:} At 8 threads, efficiency drops to 17--18\%, indicating serial bottlenecks (migration, reduction, I/O) consume $\approx$80\% of time.
	
	\item \textbf{Topology Comparison:} Star performs slightly better at 2 threads (1.03 speedup vs 0.81); Ring catches up at 4+ threads. Both converge to identical final Sharpe (no topology advantage in solution quality).
\end{enumerate}

\subsection{Convergence Analysis}

\begin{longtable}[htp]{|p{2cm}|p{2cm}|p{2cm}|p{2cm}|p{2cm}|}
\caption{Convergence Milestones (Ring Topology, 1 Thread)}
\label{Table:Convergence}
\hline
\textbf{Generation} & \textbf{Mean Best Sharpe} & \textbf{Std Best Sharpe} & \textbf{Improvement} & \textbf{Status} \\
\hline
1 & $0.043452$ & $0.001817$ & -- & Initial \\
\hline
10 & $0.052481$ & $0.001101$ & $+20.9\%$ & Early phase \\
\hline
50 & $0.066055$ & $0.000258$ & $+52.1\%$ & Mid phase \\
\hline
100 & $0.067063$ & $0.000014$ & $+54.5\%$ & Near plateau \\
\hline
150 & $0.067188$ & $0.000009$ & $+54.7\%$ & Convergence \\
\hline
200 & $0.067242$ & $0.000010$ & $+54.8\%$ & Converged \\
\hline
\end{longtable}

The algorithm exhibits rapid early convergence (gens 1--50) followed by slower refinement (gens 50--200). This is typical of GA behavior: quick discovery of good regions, then fine-tuning.

\subsection{Parallel Efficiency Analysis}

Observed speedup vs.~theoretical upper bound (Amdahl's Law):

\begin{equation}
\text{Speedup} = \frac{1}{f_{\text{serial}} + (1 - f_{\text{serial}}) / p}
\label{eq:amdahl}
\end{equation}

\noindent where $f_{\text{serial}}$ is the serial fraction and $p$ is the number of threads.

Estimated serial fraction from Ring topology:
\begin{equation}
f_{\text{serial}} \approx 1 - \frac{\text{ideal speedup at 8 threads}}{8} \approx 1 - \frac{1.40}{8} \approx 0.825
\label{eq:serial_frac}
\end{equation}

This indicates \textbf{82.5\% of time is spent in serial phases} (migration, I/O, reductions), limiting speedup to $\approx 1.4 \times$ at 8 threads even in the ideal case. Memory bandwidth contention and thread scheduling overhead further reduce observed speedup.

\subsection{Sharpe Ratio Interpretation}

Daily Sharpe of $0.0672$:
\begin{equation}
\text{Annualized Sharpe} = 0.0672 \times \sqrt{252} \approx 1.064
\label{eq:annual_sharpe}
\end{equation}

This is \textbf{reasonable} for a long-only portfolio over the 2010--2021 period, which included:
\begin{itemize}
	\item Post-financial crisis recovery (2010--2012)
	\item 2014--2015 market corrections
	\item 2020 COVID crash
\end{itemize}

The GA successfully navigated these volatile periods and found a diversified portfolio with attractive risk-adjusted returns.

%==============================================================================
\newpage
\section{Bottleneck Analysis and Scaling Limitations}

\subsection{Serial Phases}

\begin{itemize}
	\item \textbf{Migration (Ring):} Island $i$ sends top individuals to island $(i+1)$. O($m_{\text{mig}} \times k$) complexity, runs serially after each island generation.
	
	\item \textbf{Migration (Star):} Hub collects all migrants, sorts globally, broadcasts. O($k \times m_{\text{mig}} \log(k \times m_{\text{mig}})$) for sorting, higher overhead than Ring.
	
	\item \textbf{Global best tracking:} Reduction across islands to find maximum Sharpe. O($k$) serial work.
	
	\item \textbf{I/O:} CSV write for convergence and performance data. Amortized once per run (negligible per generation).
\end{itemize}

\subsection{Memory Bandwidth}

At 8 threads, fitness evaluation for $P = 80$ individuals may become memory-bound if the covariance matrix ($20 \times 20 = 400$ floats) is repeatedly loaded from memory. Modern CPUs can sustain $\approx 60$ GB/s; with 8 threads, this may be exhausted.

\subsection{Recommendations for Further Scaling}

\begin{itemize}
	\item \textbf{Increase population or islands:} Larger $P$ and $k$ dilute serial fraction (strong scaling study).
	
	\item \textbf{Use MPI:} For distributed memory on compute clusters; eliminates shared-memory thread limits.
	
	\item \textbf{Reduce migration frequency:} Fewer migrations = less serialization (trade-off: slower convergence).
	
	\item \textbf{Lock-free data structures:} Enable finer parallelism inside GA steps (selection, crossover) at implementation complexity cost.
\end{itemize}

%==============================================================================
\newpage
\section{Conclusion}

This project successfully implemented an Island Model Genetic Algorithm with two-level OpenMP parallelism (task + data) to solve a realistic portfolio optimization problem on NIFTY 50 historical data. Key contributions include:

\begin{itemize}
	\item \textbf{Algorithm:} IMGA with Sharpe ratio fitness on mean-variance portfolios, achieving daily Sharpe $\approx 0.0672$ (annualized $\approx 1.064$).
	
	\item \textbf{Parallelization:} Demonstrated task parallelism (island-level) and data parallelism (fitness evaluation), with measured speedup up to $1.40 \times$ at 8 threads.
	
	\item \textbf{Topology Comparison:} Ring vs.~Star migrations show similar final performance; Ring slightly more efficient due to lower communication overhead.
	
	\item \textbf{Reproducibility:} 5 trials per configuration with statistical averaging confirm robustness; std.~deviation in final Sharpe $< 10^{-5}$.
	
	\item \textbf{Scaling Analysis:} Serial phases (migration, I/O, reduction) limit speedup to $\approx 1.4 \times$; Amdahl's Law explains observed inefficiency at 8 threads.
\end{itemize}

The project demonstrates that OpenMP on shared-memory systems is effective for moderate-scale portfolio optimization, with clear bottlenecks identified for future optimization (MPI, lock-free structures, larger datasets).

%==============================================================================
\newpage
\section{Future Work}

\begin{enumerate}
	\item \textbf{MPI Extension:} Implement true distributed Island GA on compute clusters for 100+ cores.
	
	\item \textbf{Weak Scaling Study:} Increase population and generations with thread count to assess weak scaling efficiency.
	
	\item \textbf{VTune / Perf Profiling:} Detailed breakdown of where time is spent within each phase.
	
	\item \textbf{Larger Datasets:} Extend to 50--100 stocks and multi-decade windows to stress-test algorithm and parallelism.
	
	\item \textbf{Advanced Migration:} Implement adaptive migration frequency based on convergence stagnation.
\end{enumerate}

%==============================================================================
\newpage
\addcontentsline{toc}{section}{References}
\bibliographystyle{IEEEtran}
\bibliography{mybib}

\end{document}
